
\section*{Wnioski}
\addcontentsline{toc}{section}{Wnioski}

Na podstawie przeprowadzonych prac nad systemem RiskScoop AI sformułowano następujące wnioski:

\begin{enumerate}
    \item \textbf{skuteczność agentów AI} -- duże modele językowe upraszczają interakcję użytkownika z~systemem analizy geoprzestrzennej -- użytkownik nie musi znać składni GIS,

    \item \textbf{wartość otwartych danych} -- integracja z~Overture Maps Foundation i~Copernicus GloFAS pokazuje, że otwarte dane wystarczają do budowy użytecznych narzędzi analitycznych,

    \item \textbf{architektura AgentOS} -- wzorzec AgentOS z~frameworka Agno upraszcza implementację -- automatyczne generowanie API eliminuje ręczne definiowanie endpointów,

    \item \textbf{model ryzyka} -- klasyfikacja oparta na 5~klasach głębokości powodzi z~danych GloFAS okazała się skuteczna w~identyfikacji zagrożonej infrastruktury.
\end{enumerate}

\subsection*{Ograniczenia pracy}

\begin{itemize}
    \item zależność od zewnętrznych API (Google Gemini, Copernicus CDS, Overture Maps),
    \item dane GloFAS reprezentują scenariusze modelowe, nie rzeczywiste bieżące zagrożenie,
    \item brak walidacji modelu z~danymi historycznymi o~rzeczywistych stratach powodziowych.
\end{itemize}

\subsection*{Kierunki dalszych prac}

\begin{enumerate}
    \item integracja z~polskim systemem ISOK i~danymi IMGW,
    \item rozbudowa modelu o~współczynniki podatności infrastruktury,
    \item generowanie raportów PDF dla potrzeb zarządzania kryzysowego,
    \item powiadomienia o~zmianie poziomu ryzyka w~czasie rzeczywistym.
\end{enumerate}

\section*{Podsumowanie}
\addcontentsline{toc}{section}{Podsumowanie}

Praca przedstawiła projekt i~implementację systemu RiskScoop AI -- narzędzia do analizy ryzyka powodziowego z~wykorzystaniem agenta AI.

\subsection*{Realizacja celów pracy}

Wszystkie cele zdefiniowane we wstępie zostały zrealizowane:

\begin{enumerate}
    \item \textbf{cel główny} -- opracowano działający system analizy ryzyka powodziowego z~interfejsem w~języku naturalnym,

    \item \textbf{cele szczegółowe}:
    \begin{itemize}
        \item zaprojektowano architekturę agenta AI z~6~narzędziami,
        \item zintegrowano dane Overture Maps, Copernicus GloFAS i~OpenStreetMap,
        \item zaimplementowano model klasyfikacji ryzyka,
        \item stworzono responsywny interfejs z~mapą Mapbox.
    \end{itemize}
\end{enumerate}

\subsection*{Weryfikacja kryteriów sukcesu}

Wszystkie kryteria sukcesu zdefiniowane w~rozdziale ,,Cel i~zakres pracy'' zostały spełnione:

\begin{itemize}
    \item interpretacja języka naturalnego: \textbf{80\%} (wymagane $\geq$80\%),
    \item czas odpowiedzi: \textbf{90\% zapytań $\leq$30s} (wymagane $\leq$30s),
    \item identyfikacja obiektów: zgodna z~danymi Overture Maps,
    \item kwantyfikacja ryzyka: oparta na danych GloFAS,
    \item interfejs użytkownika: intuicyjny, z~filtrowaniem i~wizualizacją.
\end{itemize}

\subsection*{Główne osiągnięcia}

\begin{enumerate}
    \item agent AI z~AgentOS i~modelem Gemini 3 Pro,
    \item 6 narzędzi do analizy geoprzestrzennej,
    \item algorytm \texttt{assemble\_rings} do rekonstrukcji granic z~OSM,
    \item integracja trzech źródeł danych geoprzestrzennych,
    \item interfejs użytkownika z~mapą interaktywną.
\end{enumerate}

\subsection*{Praktyczne zastosowania}

System może znaleźć zastosowanie w:

\begin{itemize}
    \item zarządzaniu kryzysowym -- ocena zagrożeń dla infrastruktury,
    \item planowaniu przestrzennym -- identyfikacja obszarów ryzyka,
    \item ubezpieczeniach -- szacowanie ekspozycji na ryzyko powodziowe.
\end{itemize}

\vspace{1cm}

Zrealizowany prototyp pokazuje, że agenci AI w~połączeniu z~otwartymi danymi geoprzestrzennymi pozwalają budować narzędzia analityczne dostępne dla użytkowników bez wiedzy technicznej.

